\documentclass{article}
\usepackage{amsmath,amssymb,amsthm}
\usepackage[margin=1in]{geometry}

\title{Reading Notes: Automata Theory Meets Barrier Certificates\\
Temporal Logic Verification of Nonlinear Systems}
\author{Wongpiromsarn, Topcu, Lamperski (IEEE TAC 2016)}
\date{Paper Read: 2025-12-28}

\newtheorem{definition}{Definition}
\newtheorem{lemma}{Lemma}
\newtheorem{theorem}{Theorem}
\newtheorem{corollary}{Corollary}

\begin{document}

\maketitle

\section{Paper Metadata}

\begin{itemize}
    \item \textbf{Title}: Automata Theory Meets Barrier Certificates: Temporal Logic Verification of Nonlinear Systems
    \item \textbf{Authors}: Tichakorn Wongpiromsarn, Ufuk Topcu, Andrew Lamperski
    \item \textbf{Publication}: IEEE Transactions on Automatic Control, Vol. 61, No. 11, November 2016
    \item \textbf{Pages}: 3344-3355
    \item \textbf{DOI}: 10.1109/TAC.2015.2511722
    \item \textbf{Received}: January 8, 2015
    \item \textbf{Accepted}: December 13, 2015
\end{itemize}

\section{Core Contribution and Approach}

This paper proposes a \textbf{sound but incomplete} method for computational verification of LTL specifications against dynamical systems over continuous state spaces by merging:
\begin{enumerate}
    \item Automata-based model checking
    \item Barrier certificates from control theory
    \item Sum-of-squares (SOS) optimization for certificate synthesis
\end{enumerate}

\textbf{Key Innovation}: The automaton representation of $\neg\varphi$ (negation of LTL formula) guides decomposition of verification task into a finite collection of simpler constraints over continuous state space. These constraints are then proved using barrier certificates, avoiding explicit finite-state abstractions.

\section{System Model}

Consider a dynamical system $\mathfrak{D}$ with state $x \in \mathcal{X} \subseteq \mathbb{R}^n$ evolving according to:
\begin{equation}
\dot{x}(t) = f(x(t))
\end{equation}
where $f$ is Lipschitz continuous to ensure piecewise continuously differentiable solutions.

\subsection{Atomic Propositions and Regions}

\begin{itemize}
    \item Set of atomic propositions: $\Pi = \{p_0, p_1, \ldots, p_N\}$
    \item Each $p_i$ corresponds to a region $\llbracket p_i \rrbracket = X_i \subseteq \mathcal{X}$
    \item For each $a \in 2^\Pi$, define:
\end{itemize}
\begin{equation}
\llbracket a \rrbracket = \begin{cases}
\mathcal{X} \setminus \bigcup_{p \in \Pi} \llbracket p \rrbracket & \text{if } a = \emptyset \\
\bigcap_{p \in a} \llbracket p \rrbracket \setminus \bigcup_{p \in \Pi \setminus a} \llbracket p \rrbracket & \text{otherwise}
\end{cases}
\end{equation}

\section{Barrier Certificate Foundation}

\begin{lemma}[Basic Barrier Certificate]
Let $\mathcal{Y}, \mathcal{Y}_0, \mathcal{Y}_1 \subseteq \mathcal{X}$. Suppose there exists a differentiable function $B: \mathcal{X} \to \mathbb{R}$ satisfying:
\begin{align}
B(x) &\leq 0 \quad \forall x \in \mathcal{Y}_0 \label{eq:bc1}\\
B(x) &> 0 \quad \forall x \in \mathcal{Y}_1 \label{eq:bc2}\\
\frac{\partial B}{\partial x}(x) f(x) &\leq 0 \quad \forall x \in \mathcal{Y} \setminus \mathcal{Y}_1 \label{eq:bc3}
\end{align}
Then, there exists no trajectory $x$ of $\mathfrak{D}$ such that $x(0) \in \mathcal{Y}_0$, $x(T) \in \mathcal{Y}_1$ and $x(t) \in \mathcal{Y}$ for all $t \in [0,T]$.
\end{lemma}

\textbf{Interpretation}: $B$ is a barrier that:
\begin{itemize}
    \item Is non-positive at initial states $\mathcal{Y}_0$
    \item Is positive at target states $\mathcal{Y}_1$
    \item Is non-increasing along trajectories in $\mathcal{Y} \setminus \mathcal{Y}_1$
\end{itemize}
Thus, no trajectory starting in $\mathcal{Y}_0$ can reach $\mathcal{Y}_1$ while staying in $\mathcal{Y}$.

\section{Trace Semantics}

\begin{definition}[Trace]
An infinite sequence $\sigma_x = a_0 a_1 a_2 \ldots$ where $a_i \in 2^\Pi$ is a \textbf{trace} of trajectory $x: \mathbb{R}_{\geq 0} \to \mathcal{X}$ if there exists a sequence $t_0 t_1 t_2 \ldots$ with $t_0 = 0$, $t_k \to \infty$ such that:
\begin{enumerate}
    \item $t_i < t_{i+1}$
    \item $x(t_i) \in \llbracket a_i \rrbracket$
    \item If $a_i \neq a_{i+1}$, then $\exists t_i' \in [t_i, t_{i+1}]$ such that:
    \begin{itemize}
        \item $x(t) \in \llbracket a_i \rrbracket$ for all $t \in (t_i, t_i')$
        \item $x(t) \in \llbracket a_{i+1} \rrbracket$ for all $t \in (t_i', t_{i+1})$
        \item Either $x(t_i') \in \llbracket a_i \rrbracket$ or $x(t_i') \in \llbracket a_{i+1} \rrbracket$
    \end{itemize}
\end{enumerate}
\end{definition}

\textbf{Key Point}: A trace captures the sequence of regions (characterized by atomic propositions) visited along the trajectory. The trace may not be unique for a given trajectory.

\subsection{Satisfaction by System}

\begin{definition}[System Satisfies LTL]
System $\mathfrak{D}$ satisfies LTL formula $\varphi$ if:
\begin{equation}
\text{Trace}(\mathfrak{D}) \subseteq \text{Words}(\varphi)
\end{equation}
where $\text{Trace}(\mathfrak{D})$ is the set of all traces of all trajectories of $\mathfrak{D}$, and $\text{Words}(\varphi) = \{\sigma \in (2^\Pi)^\omega \mid \sigma \models \varphi\}$.
\end{definition}

\section{The State-Triplet Method}

\subsection{Core Idea}

Given NBA $\mathcal{A}_{\neg\varphi} = (Q, 2^\Pi, \delta, Q_0, F)$ for $\neg\varphi$:

\begin{enumerate}
    \item Enumerate all \textbf{simple paths} from initial states $Q_0$ to accepting states $F$
    \item For each simple path $q_0 \to q_1 \to \cdots \to q_m$ where $q_m \in F$, use barrier certificates to prove:
    \begin{itemize}
        \item No trajectory can follow the sequence of regions corresponding to this path
    \end{itemize}
    \item Block each simple path individually using one barrier certificate per path
\end{enumerate}

\subsection{Conservatism of State-Triplet Method}

\textbf{Major Limitation}: The method requires:
\begin{itemize}
    \item Unrolling the automaton to enumerate all simple paths
    \item Blocking \textbf{all} simple paths to accepting states
    \item Cannot handle cases where accepting states are reachable but accepting transitions occur finitely often
\end{itemize}

\textbf{Example Where State-Triplet Fails}: Consider verifying $\square \neg a$ against NBA for $\Diamond a$. The accepting state $q_0$ has self-loop. If $q_0$ is reachable, state-triplet method fails, even if the system only visits states satisfying $a$ finitely many times.

\section{This Paper's Method: Automata-Based Decomposition}

\subsection{Key Construction}

\textbf{Step 1}: Construct set $\Omega \subseteq (2^\Pi)^*$ of finite strings such that:
\begin{lemma}
For any word $\sigma \in (2^\Pi)^\omega$, if $\sigma \not\models \varphi$, then there exists a substring $\omega \in \Omega$ of $\sigma$.
\end{lemma}

\textbf{Step 2}: For each $\omega \in \Omega$, prove that $\omega$ cannot be a substring of any word in $\text{Trace}(\mathfrak{D})$.

\subsection{Construction of $\Omega$}

Given NBA $\mathcal{A}_{\neg\varphi} = (Q, 2^\Pi, \delta, Q_0, F)$:

\begin{enumerate}
    \item Define $\mathcal{R}^{\text{acc}} = \{\pi^p \pi^c \mid \pi^p \in \mathcal{R}(q_0, q), \pi^c \in \mathcal{R}(q', q), q_0 \in Q_0, q \in F, (q, a, q') \in \delta \text{ for some } a\}$

    where:
    \begin{itemize}
        \item $\pi^p$: path from initial state to accepting state $q$
        \item $\pi^c$: accepting cycle starting from $q$
    \end{itemize}

    \item Define:
    \begin{equation}
    \Omega = \bigcup_{\pi \in \mathcal{R}^{\text{acc}}} \mathcal{ST}(\pi)
    \end{equation}
    where $\mathcal{ST}(\pi)$ is the set of all strings generated by run fragment $\pi$.
\end{enumerate}

\textbf{Problem}: $\mathcal{R}^{\text{acc}}$ and hence $\Omega$ are typically \textbf{infinite}!

\section{Representative Run Fragments}

To handle infinite $\Omega$, construct finite representative sets:

\subsection{Path Sets Without Repeated Edges}

For each accepting state $q \in F$, compute:
\begin{itemize}
    \item $\mathcal{P}^{\text{cyc}}(q)$: Paths from $q$ to $q$ with no repeated edges and no consecutive state repetitions
    \item $\mathcal{P}^{\text{path}}(q)$: Paths from initial states to $q$ with no repeated edges and no consecutive state repetitions
\end{itemize}

Both sets are \textbf{finite} (can be computed via DFS - Algorithm 1 in paper).

\subsection{Length-3 Subpaths}

For finite path $\pi = q_0 q_1 \ldots q_m$, define:
\begin{equation}
\mathcal{PF}^3(\pi) = \{q_i q_{i+1} q_{i+2} \mid 0 \leq i \leq m-2\}
\end{equation}

\textbf{Why length 3?} Any path can be "extended" from its length-3 subpaths by inserting consecutive state repetitions. Length $< 3$ is not useful; length $> 3$ increases complexity.

\section{Main Verification Result}

\begin{theorem}[Proposition 1 in paper]
Suppose for each $q \in F$, either condition (1) or (2) holds:

\textbf{Condition (1)} [Invalidate Cycles]: For each $p \in \mathcal{P}^{\text{cyc}}(q)$, there exists $\pi = q_0 q_1 q_2 \in \mathcal{PF}^3(p)$ such that:
\begin{itemize}
    \item[(a)] All strings $a_0 a_1 \in \mathcal{ST}(\pi)$ cannot be substrings of any word in $\text{Trace}(\mathfrak{D})$
    \item[(b)] If $(q_1, q_1) \in E^G$, then all strings $a_0 \tilde{a}_0 \ldots \tilde{a}_k a_1 \in \mathcal{ST}(q_0 q_1^+ q_2)$ cannot be substrings of any word in $\text{Trace}(\mathfrak{D})$
\end{itemize}

\textbf{Condition (2)} [Invalidate Paths]: For each $p \in \mathcal{P}^{\text{path}}(q)$, there exists $\pi = q_0 q_1 q_2 \in \mathcal{PF}^3(p)$ such that conditions (a) and (b) above hold.

Then $\mathfrak{D} \models \varphi$.
\end{theorem}

\section{Barrier Certificate Conditions}

\subsection{For Condition (a): Invalidating $a_0 a_1$}

\begin{corollary}[Corollary 1]
Consider $\Sigma_0, \Sigma_1 \subseteq 2^\Pi$ and $\tilde{\Omega} = \{a_0 a_1 \mid a_0 \in \Sigma_0, a_1 \in \Sigma_1\}$.

Let $\mathcal{Y}_0 = \bigcup_{a \in \Sigma_0} \llbracket a \rrbracket$, $\mathcal{Y}_1 = \bigcup_{a \in \Sigma_1} \llbracket a \rrbracket$, $\mathcal{Y} = \mathcal{Y}_0 \cup \mathcal{Y}_1$.

If there exists differentiable $B: \mathcal{X} \to \mathbb{R}$ satisfying:
\begin{align}
B(x) &\leq 0 \quad \forall x \in \mathcal{Y}_0\\
B(x) &> 0 \quad \forall x \in \mathcal{Y}_1\\
\frac{\partial B}{\partial x}(x) f(x) &\leq 0 \quad \forall x \in \mathcal{Y} \setminus \mathcal{Y}_1
\end{align}
Then no string in $\tilde{\Omega}$ can be a substring of any word in $\text{Trace}(\mathfrak{D})$.
\end{corollary}

\subsection{For Condition (b): Invalidating $a_0 \tilde{a}_0 \ldots \tilde{a}_k a_1$}

\begin{corollary}[Corollary 2]
Consider $\Sigma_0, \Sigma_1, \tilde{\Sigma} \subseteq 2^\Pi$ and:
\begin{equation}
\tilde{\Omega} = \{a_0 \tilde{a}_0 \ldots \tilde{a}_k a_1 \mid k \in \mathbb{N}, a_0 \in \Sigma_0, \tilde{a}_0,\ldots,\tilde{a}_k \in \tilde{\Sigma}, a_1 \in \Sigma_1\}
\end{equation}

Let $\mathcal{Y}_0 = \bigcup_{a \in \Sigma_0} \llbracket a \rrbracket$, $\mathcal{Y}_1 = \bigcup_{a \in \Sigma_1} \llbracket a \rrbracket$, $\tilde{\mathcal{Y}} = \bigcup_{a \in \tilde{\Sigma}} \llbracket a \rrbracket$, $\mathcal{Y} = \mathcal{Y}_0 \cup \mathcal{Y}_1 \cup \tilde{\mathcal{Y}}$.

If there exists differentiable $B: \mathcal{X} \to \mathbb{R}$ satisfying barrier certificate conditions (same as Corollary 1), then no string in $\tilde{\Omega}$ can be a substring of any word in $\text{Trace}(\mathfrak{D})$.
\end{corollary}

\section{SOS-Based Synthesis}

\begin{lemma}[Lemma 7 - SOS Sufficient Conditions]
Let $\mathcal{Y}, \mathcal{Y}_0, \mathcal{Y}_1 \subseteq \mathcal{X}$ be defined by:
\begin{align}
\mathcal{Y}_0 &= \{x \in \mathbb{R}^n \mid g_0(x) \geq 0\}\\
\mathcal{Y}_1 &= \{x \in \mathbb{R}^n \mid g_1(x) \geq 0\}\\
\mathcal{Y} &= \{x \in \mathbb{R}^n \mid g(x) \geq 0\}
\end{align}

Suppose there exist polynomial $B$, constant $\epsilon > 0$, and sum-of-squares polynomials $s_0, s_1, s_2, s_3$ such that the following are SOS:
\begin{align}
-B(x) - s_0(x) g_0(x) &\text{ is SOS} \label{eq:sos1}\\
B(x) - \epsilon - s_1(x) g_1(x) &\text{ is SOS} \label{eq:sos2}\\
-\frac{\partial B}{\partial x}(x) f(x) - s_2(x) g(x) + s_3(x) g_1(x) &\text{ is SOS} \label{eq:sos3}
\end{align}
Then $B$ satisfies barrier certificate conditions.
\end{lemma}

\textbf{Computational Method}: The SOS conditions can be reformulated as semidefinite programs (SDPs) and solved using SOSTOOLS or similar tools.

\section{Comparison: This Paper vs. My Work}

\begin{center}
\begin{tabular}{|p{4cm}|p{5cm}|p{5cm}|}
\hline
\textbf{Aspect} & \textbf{Wongpiromsarn 2015} & \textbf{My Work} \\
\hline
\textbf{Certificate Type} & Barrier certificates on product system $\mathcal{X} \times Q$ &
\begin{itemize}
\item Safety: $B_k(x,q)$ on $\mathcal{X} \times Q$
\item Persistence: $B_{k,l}(x)$ on $\mathcal{X}$
\end{itemize} \\
\hline
\textbf{Verification Approach} & Block simple paths using path enumeration &
\begin{itemize}
\item Safety: Block reachability to $q_k \in Acc^{\text{start}}$
\item Persistence: Prove finite occurrence via ranking function
\end{itemize} \\
\hline
\textbf{Accepting Conditions} & NBA with state-based acceptance $F \subseteq Q$ & NBA with transition-based acceptance $Acc \subseteq \delta$ \\
\hline
\textbf{Search Space} & $\mathcal{X} \times Q$ for each barrier certificate & At most $\mathcal{X} \times Q \times Q$ total \\
\hline
\textbf{Handling Persistence} & Cannot verify if accepting states are reachable but visited finitely often & Can verify via ranking function even if accepting transitions occur finitely \\
\hline
\textbf{Number of Certificates} & One per simple path (potentially many) &
\begin{itemize}
\item One safety cert per $q_k \in Acc^{\text{start}}$
\item OR persistence certs for all $(k,l) \in Acc^{\text{pair}}$ with start $q_k$
\end{itemize} \\
\hline
\textbf{Conservatism} & High: must block all simple paths & Lower: allows reachability with finite occurrence \\
\hline
\textbf{Completeness} & Incomplete (sound only) & Incomplete (sound only) \\
\hline
\end{tabular}
\end{center}

\section{Key Insights for My Paper}

\subsection{What to Emphasize}

\begin{enumerate}
    \item \textbf{Reduced Conservatism}: Our method does NOT require blocking all paths to accepting states. We allow:
    \begin{itemize}
        \item Accepting states to be reachable (state-triplet cannot handle this)
        \item Accepting transitions to occur, as long as finitely often (persistence certificate)
    \end{itemize}

    \item \textbf{Smaller Search Space}:
    \begin{itemize}
        \item Wongpiromsarn: Each barrier certificate on $\mathcal{X} \times Q$, need one per simple path
        \item Our safety certificate: $B_k(x,q)$ on $\mathcal{X} \times Q$, one per $q_k \in Acc^{\text{start}}$
        \item Our persistence certificate: $B_{k,l}(x)$ on $\mathcal{X}$, total search space $\leq \mathcal{X} \times Q \times Q$
    \end{itemize}

    \item \textbf{No Path Enumeration}: We do not need to enumerate simple paths. We work directly with accepting transition start states $Acc^{\text{start}}$ and pairs $Acc^{\text{pair}}$.

    \item \textbf{Transition-Based NBA}: Using transition-based acceptance conditions is more efficient for LTL verification.

    \item \textbf{Complementarity}: Safety and persistence certificates are complementary:
    \begin{itemize}
        \item Safety blocks reachability
        \item Persistence allows reachability but proves finite occurrence via ranking
    \end{itemize}
\end{enumerate}

\subsection{Example Demonstrating Advantage}

Consider verifying $\square \neg a$ with NBA for $\Diamond a$ (Figure 1 in my paper):
\begin{itemize}
    \item NBA has accepting self-loop $(q_0, \{a\}, q_0)$ at $q_0$
    \item \textbf{Wongpiromsarn method}: Needs to prove $q_0$ unreachable (conservative)
    \item \textbf{My method}:
    \begin{itemize}
        \item Try safety certificate for $q_0$ unreachable
        \item If fails, try persistence certificate proving accepting transitions occur finitely often
        \item Persistence certificate fails (correctly) because accepting transition occurs at every step when in $\{a\}$ region
        \item But safety certificate succeeds if system never reaches states satisfying $a$
    \end{itemize}
\end{itemize}

\section{State-Triplet Method Definition}

Based on this paper (Section 3.4.1, page 3350), the \textbf{state-triplet approach}:

\begin{quote}
``uses barrier certificates to separate safe regions $\mathcal{X}_{\text{safe}}$ from unsafe regions $\mathcal{X}_u$. It blocks all simple paths from $q_0$ to accepting states in the NBA for $\neg\varphi$. Finding certificates for all such paths ensures $\mathfrak{L}(\mathfrak{S})$ never visits accepting states, verifying $\mathfrak{S} \models_\mathfrak{L} \varphi$.''
\end{quote}

\textbf{Key characteristics}:
\begin{itemize}
    \item Enumerate all simple paths from initial states to accepting states
    \item Use one barrier certificate per simple path (state triplet)
    \item Each certificate proves that path cannot be followed
    \item Conservative: cannot handle cases where accepting states are reachable
\end{itemize}

\section{Summary and Recommendation}

\subsection{How to Use This Paper in Related Work}

When writing our paper's related work section:

\begin{enumerate}
    \item \textbf{Cite as foundational work}: ``Wongpiromsarn et al.~\cite{wongpiromsarn2015automata} developed a state-triplet method that conservatively identifies barrier certificates to obstruct transitions between automaton edges, thereby preventing the system from entering accepting states.''

    \item \textbf{Highlight limitations}:
    \begin{itemize}
        \item Requires unrolling automaton and blocking all simple paths
        \item Cannot verify systems where accepting states are reachable but transitions occur finitely often
        \item Each path requires separate barrier certificate
    \end{itemize}

    \item \textbf{Position our contribution}:
    \begin{itemize}
        \item We reduce conservatism by allowing accepting states to be reachable
        \item We use transition-based NBA for more efficient verification
        \item Our persistence certificates use ranking functions to prove finite occurrence
        \item Our search space is smaller: $\mathcal{X} \times Q \times Q$ vs. multiple $\mathcal{X} \times Q$ certificates
    \end{itemize}
\end{enumerate}

\subsection{Citation Correctness}

\textbf{Verified}: This paper is \textbf{directly relevant} to our work:
\begin{itemize}
    \item Uses barrier certificates for LTL verification ✓
    \item Works with NBA (state-based acceptance) ✓
    \item Addresses temporal logic verification of dynamical systems ✓
    \item Is the foundational work for state-triplet method ✓
\end{itemize}

\textbf{Recommendation}: This citation is \textbf{essential and correct}. It should be kept and cited prominently in related work.

\end{document}
